\documentclass[12pt]{article}
\usepackage[T1]{fontenc}
\usepackage[utf8]{inputenc}
\usepackage[brazil]{babel}
\usepackage[margin=1in]{geometry}
\usepackage{ragged2e}
\usepackage[normalem]{ulem}
\setlength{\parindent}{0pt}
\setlength{\parskip}{0.8\baselineskip}
\begin{document}
\justifying
\noindent Verifica-se que a questão objeto da controvérsia jurídica suscitada no recurso manejado pela parte recorrente esteve afetada no representativo da controvérsia - \textbf{Tema n. }\textbf{246} da TNU - (PU no processo n. PEDILEF 0500881-37.2018.4.05.8204/PB),  afetado em 12/12/2019 foi julgado (trânsito em julgado em 29/01/2021), por meio do qual foi elucidada a seguinte questão:

\noindent \textit{"A partir da regra constante do art. 60, §9.º, da Lei n.º 8.213/91, saber se, para fins de fixação da DCB do auxílio-doença concedido judicialmente, o prazo de recuperação estimado pelo perito judicial deve ser computado a partir da data de sua efetiva implantação ou da data da perícia judicial."}

\noindent Restou firmada a seguinte tese:

\noindent \textit{"I - Quando a decisão judicial adotar a estimativa de prazo de recuperação da capacidade prevista na perícia, o termo inicial é a data da realização do exame, sem prejuízo do disposto no art. 479 do CPC, devendo ser garantido prazo mínimo de 30 dias, desde a implantação, para viabilizar o pedido administrativo de prorrogação. II - quando o ato de concessão (administrativa ou judicial) não indicar o tempo de recuperação da capacidade, o prazo de 120 dias, previsto no § 9º, do art. 60 da Lei 8.213/91, deve ser contado a partir da data da efetiva implantação ou restabelecimento do benefício no sistema de gestão de benefícios da autarquia." (Tema 246 TNU).}

\noindent Com efeito, depreende-se da leitura do voto condutor do acórdão que o mesmo se encontra em conformidade com o decidido no  tema pela TNU.

\noindent A proposito, consta do acordao hostilizado: (...) Adicional de Periculosidade, Sistema Remuneratório e Benefícios, Servidor Público Civil, DIREITO ADMINISTRATIVO E OUTRAS MATÉRIAS DE DIREITO PÚBLICO

\noindent Nessas condições, o conteúdo do Acórdão recorrido é consentâneo com o entendimento da TNU, firmado em julgamento de recurso representativo de controvérsia, o que atrai a incidência da regra prevista pelo art. 14, III, \textit{b,} da Resolução CJF n. 586/2019:

\noindent \textit{\textbf{Art. 14.}}\textit{ Decorrido o prazo para contrarrazões, os autos serão conclusos ao magistrado responsável pelo exame preliminar de admissibilidade, que deverá, de forma sucessiva: (...) }\textit{\textbf{III }}\textit{- negar seguimento a pedido de uniformização de interpretação de lei federal interposto contra acórdão que esteja em conformidade com entendimento consolidado: }\textit{\textbf{b)}}\textit{ em }\uline{\textit{\textbf{recurso representativo de controvérsia pela Turma Nacional de Uniformização}}}\textit{ ou em pedido de uniformização de interpretação de lei dirigido ao Superior Tribunal de Justiça;(...).}

\noindent Posto isso, com arrimo no \uline{\textbf{art. 14, III, }}\uline{\textit{\textbf{b}}}\textit{,} da Resolução CJF n. 586/2019 c/c art. 1.030, I, do CPC, \textbf{NEGO SEGUIMENTO }\textbf{a}o\textbf{ PU}\textbf{.}

\noindent Intimem-se. Aguarde-se o decurso do prazo recursal. Após, certifique-se o trânsito em julgado e devolva-se ao juízo de origem.
\end{document}
